\section*{零基础读证明：随手查阅}
\addcontentsline{toc}{section}{零基础读证明：随手查阅}
\textbf{对象的层次。} $x\in A$ 表示一个元素属于集合，$A\subseteq B$ 表示 $A$ 的每个元素都在 $B$ 中。$\{x\}$ 是只含 $x$ 的集合，不是 $x$ 本身。$\varnothing$ 是空集，$\{\varnothing\}$ 则含有一个元素。$\R^N$ 中的点是有 $N$ 个坐标的有序列表。

\textbf{量词。} $\forall$ 是“对所有”，$\exists$ 是“至少存在一个”。例如“每个学生有一本书”不要求是同一本；“有一本书供每个学生使用”则要求同一本。证明极限时 $\forall\eps>0\ \exists n_0\ \forall n\ge n_0$ 的顺序不可换：先指定精度，再找统一起点，最后控制这个起点之后的全部项。

\textbf{怎样否定。} 否定“所有 $x$ 满足 $P$”得到“存在 $x$ 不满足 $P$”；否定“存在”得到“所有都不”。否定 $P\Rightarrow Q$ 是 $P$ 真且 $Q$ 假。因此要推翻一个定理，必须满足其全部假设又违反结论；拿不满足假设的例子不能当反例。

\textbf{四种基本证明。} (1) 包含 $A\subseteq B$：任取 $x\in A$，用定义得到 $x\in B$。(2) 相等 $A=B$：两个方向分别证明。(3) 等价 $P\iff Q$：证明 $P\Rightarrow Q$ 及反向。(4) 反证：把结论否定加入已知，推出矛盾。逆否命题 $\neg Q\Rightarrow\neg P$ 与 $P\Rightarrow Q$ 等价，“反过来” $Q\Rightarrow P$ 则未必成立。

\textbf{数列、子列和收敛。} 数列 $x_n$ 是第 $n$ 个对象；子列 $x_{n_k}$ 的指标必须严格递增。$x_n\to x$ 指任意小的误差最终都能保持。前面有限多项不影响极限。抽子列通常是为了用有界性和有限维 Bolzano--Weierstrass 定理获得极限；抽取之后须把原来所有相关量一起限制到同一子列。

\textbf{距离的关键估计。} $\norm{x}=(\sum_i x_i^2)^{1/2}$，$\norm{x-y}$ 是距离。三角不等式 $\norm{x-z}\le\norm{x-y}+\norm{y-z}$ 允许经中间点拆误差。两段误差各小于 $\eps/2$，总误差便小于 $\eps$；分成 $m$ 段时可各分配 $\eps/m$。这里的 $\eps$ 不是固定的小数，而是任意正精度。

\textbf{做证明题的五行草稿。} 先写完整假设；写准确目标；列出可用定义与定理；逐步注明所用条件；最后检查任意性、存在性和量词顺序。出现“显然”时问：若阅卷者要求展开，我能否补出一个集合等价链、一个不等式或一个明确构造？不知道证明的命题要标为引用，不能把待证结论循环当作理由。
\clearpage
